\documentclass[12pt]{amsart}
%\usepackage[margin=1in]{geometry}
\pagestyle{plain}

\usepackage{amsmath,amssymb}
\usepackage{enumerate}

\usepackage[shortlabels]{enumitem}
% Set spacing for all enumerate environments
\setlist[enumerate]{itemsep=6pt, topsep=6pt, parsep=0pt}

\usepackage{graphicx}

\usepackage[left=0.75in,right=1in,bottom=0.9in]{geometry}  
% \usepackage{mathptmx}      % use Times fonts if available on your TeX system
\usepackage[T1]{fontenc}

\usepackage[parfill]{parskip}
\renewcommand{\baselinestretch}{1}

% Macro definitions
\newcommand{\norm}[1]{\left\|#1\right\|}

%%%%%%%%%%%%%%%%%%%%%%%%%%%%%%%%%
% SECTION DIVIDERS
%%%%%%%%%%%%%%%%%%%%%%%%%%%%%%%%%
\newcommand{\divider}{
  \vspace{-1em} % Add space above
  \noindent
  \raisebox{-0.15ex}{\textbullet}%
  \hspace{0.5em}%
  \rule[0.5ex]{0.95\textwidth}{1.0pt}%
  \hspace{0.5em}%
  \raisebox{-0.15ex}{\textbullet}%
  \vspace{0em} % Adjust space below
}

\title{ESE 2030 QUIZZAM 5 : Fall 2025}
\author{prof-g}


\begin{document}
\maketitle

\begin{center}
    {\bf INSTRUCTIONS}
\end{center}

\begin{itemize}
\item 
No books, notes, or calculators. Use a writing utensil and logic. 

\item 
Please follow question instructions and fill in the bubble-sheet. 

\item 
Select one answer per problem.

\item 
Each problem has one best answer. 

\item 
In some problems, a wrong choice may be worth partial credit. 

\item 
Be careful to fill in the correct problem number. 

\item 
These questions are written carefully: if you detect an error, read closely and do your best.

\item 
If anything is unclear, use your best judgment. 

\item 
Cheating in any form will be dealt with severely. 
\end{itemize}

\newpage



% QUIZZAM 5
% TOTAL QUESTIONS: ~20
% COVERAGE: WEEKS 10-11


{\bf PROBLEM 1:}
% WEEK = 10
% TOPICS = SVD from geometric picture, singular value conventions, left vs right singular vectors

The figure shows how a matrix $A \in \mathbb{R}^{2 \times 2}$ transforms the unit circle to an ellipse, with the singular vectors represented.

\begin{center}
\begin{figure}[h]
\includegraphics[width=6.5in]{2030 Q5 SVD FIG.jpg}
%\caption{[left] for PROBLEM 14 ; [right] for PROBLEM 15.}
\end{figure}
\end{center}

Based on this figure, which of the following is the most likely SVD of $A = U\Sigma V^T$?

\begin{enumerate}[(A)]
\item $U = \frac{1}{\sqrt{2}}\begin{bmatrix} 1 & 1 \\ 1 & -1 \end{bmatrix}, \quad \Sigma = \begin{bmatrix} 3/2 & 0 \\ 0 & 1/2 \end{bmatrix}, \quad V = \frac{1}{\sqrt{2}}\begin{bmatrix} 1 & 1 \\ -1 & 1 \end{bmatrix}$

\item $U = \frac{1}{\sqrt{2}}\begin{bmatrix} 1 & 1 \\ -1 & 1 \end{bmatrix}, \quad \Sigma = \begin{bmatrix} 3/2 & 0 \\ 0 & 1/2 \end{bmatrix}, \quad V = \frac{1}{\sqrt{2}}\begin{bmatrix} 1 & 1 \\ 1 & -1 \end{bmatrix}$

\item $U = \frac{1}{\sqrt{2}}\begin{bmatrix} 1 & 1 \\ -1 & 1 \end{bmatrix}, \quad \Sigma = \begin{bmatrix} 1/2 & 0 \\ 0 & 3/2 \end{bmatrix}, \quad V = \frac{1}{\sqrt{2}}\begin{bmatrix} 1 & 1 \\ -1 & 1 \end{bmatrix}$

\item $U = \frac{1}{\sqrt{2}}\begin{bmatrix} 1 & 1 \\ 1 & -1 \end{bmatrix}, \quad \Sigma = \begin{bmatrix} 3/2 & 0 \\ 0 & -1/2 \end{bmatrix}, \quad V = \frac{1}{\sqrt{2}}\begin{bmatrix} 1 & 1 \\ 1 & -1 \end{bmatrix}$

\item $U = \frac{1}{\sqrt{2}}\begin{bmatrix} 1 & 1 \\ 1 & -1 \end{bmatrix}, \quad \Sigma = \begin{bmatrix} 3/2 & 0 \\ 0 & 1/2 \end{bmatrix}, \quad V = \frac{1}{\sqrt{2}}\begin{bmatrix} 1 & 1 \\ 1 & -1 \end{bmatrix}$
\end{enumerate}

% \textbf{Correct Answer:} B
% \textbf{Explanation:} Reading from the figure:
% - The right singular vectors (columns of $V$) are in the domain at $\pm 45°$: 
%   $\mathbf{v}_1 = (1,1)^T/\sqrt{2}$ and $\mathbf{v}_2 = (1,-1)^T/\sqrt{2}$
% - The semi-axis lengths give $\sigma_1 = 3/2$ (major) and $\sigma_2 = 1/2$ (minor)
% - The left singular vectors (columns of $U$) align with the ellipse axes in the codomain:
%   $\mathbf{u}_1 = (1,-1)^T/\sqrt{2}$ (at $-45°$) and $\mathbf{u}_2 = (1,1)^T/\sqrt{2}$ (at $45°$)
% 
% \textbf{Partial credit:} None
% \textbf{Trick answer:} A (swaps $U$ and $V$; confuses which singular vectors live in domain vs codomain)
% \textbf{Trick answer:} C (places smaller singular value first; violates the convention $\sigma_1 \geq \sigma_2 \geq \cdots \geq 0$)
% \textbf{Trick answer:} D (uses negative singular value $-1/2$; singular values must be non-negative by definition)
% \textbf{Trick answer:} E (uses $\sigma_i^2$ instead of $\sigma_i$; confuses eigenvalues of $A^TA$ with singular values)

\divider

%%%%%%%%%%%%%%%%%%%%%%%%%%%%%%%%%%%%%%%%%%%%%%%%%%%%%%%%%%%%%%%%%
\newpage
%%%%%%%%%%%%%%%%%%%%%%%%%%%%%%%%%%%%%%%%%%%%%%%%%%%%%%%%%%%%%%%%%


{\bf PROBLEM 2:}
% WEEK = 10
% TOPICS = Geometric foundations of SVD, sphere to ellipsoid transformation, role of A^TA

A linear transformation $A \in \mathbb{R}^{3 \times 2}$ maps the unit circle in $\mathbb{R}^2$ to an ellipse in $\mathbb{R}^3$. The matrix $A^TA$ has eigenvalues $\lambda_1 = 9$ and $\lambda_2 = 4$ with corresponding orthonormal eigenvectors $\mathbf{v}_1$ and $\mathbf{v}_2$.

Which statement correctly describes the geometric action of $A$?

\begin{enumerate}[(A)]
\item The ellipse in $\mathbb{R}^3$ has semi-axis lengths 9 and 4, aligned with the directions $\mathbf{v}_1$ and $\mathbf{v}_2$.
\item The ellipse in $\mathbb{R}^3$ has semi-axis lengths 3 and 2, and the input directions $\mathbf{v}_1$ and $\mathbf{v}_2$ are mapped to these semi-axes.
\item The ellipse lies in a 2-dimensional subspace of $\mathbb{R}^3$, with semi-axis lengths $\sqrt{9} = 3$ and $\sqrt{4} = 2$.
\item The transformation stretches the direction $\mathbf{v}_1$ by factor 9 and $\mathbf{v}_2$ by factor 4, producing an ellipse with area $36\pi$.
\item The eigenvalues of $A$ determine the semi-axis lengths, which are 9 and 4.
\end{enumerate}

% \textbf{Correct Answer:} C
% \textbf{Explanation:} The key geometric insight is that $A$ transforms the unit circle (sphere in $\mathbb{R}^2$) into an ellipse in $\mathbb{R}^3$. Since $A \in \mathbb{R}^{3 \times 2}$, the image must lie in a 2-dimensional subspace (the column space of $A$). The eigenvalues $\lambda_i$ of $A^TA$ give $\|A\mathbf{v}_i\|^2 = \lambda_i$, so the singular values $\sigma_i = \sqrt{\lambda_i}$ are the actual stretching factors: $\sigma_1 = 3$ and $\sigma_2 = 2$. These are the semi-axis lengths of the ellipse. The eigenvectors $\mathbf{v}_1, \mathbf{v}_2$ (in the domain) are the principal input directions that get stretched to form the semi-axes.
% \textbf{Partial credit:} B (recognizes semi-axis lengths are 3 and 2, but the statement about mapping is slightly imprecise—$\mathbf{v}_1, \mathbf{v}_2$ are in $\mathbb{R}^2$ while the semi-axes are in $\mathbb{R}^3$)
% \textbf{Trick answer:} A (confuses eigenvalues $\lambda_i$ with singular values $\sigma_i$; the eigenvalues give squared stretching, not the lengths themselves)
% \textbf{Trick answer:} D (uses eigenvalues as stretching factors instead of their square roots)
% \textbf{Trick answer:} E (conflates eigenvalues of $A$ with eigenvalues of $A^TA$; for rectangular $A$, eigenvalues of $A$ don't exist)

\divider


{\bf PROBLEM 3:}
% WEEKS = 10, 11 SYNTHESIS  
% TOPICS = Rank-k approximation from both perspectives, preprocessing effects

A centered data matrix $X \in \mathbb{R}^{n \times d}$ has SVD $X = U\Sigma V^T$ with singular values $\sigma_1 = 100, \sigma_2 = 50, \sigma_3 = 25, \ldots$ A researcher wants a 2-dimensional representation of the $n$ observations and considers two approaches:

\textbf{Approach 1:} Compute the $n \times 2$ matrix $U_2\Sigma_2$, where $U_2$ contains the first 2 columns of $U$ and $\Sigma_2$ is the $2 \times 2$ upper-left block of $\Sigma$.

\textbf{Approach 2:} Compute the $n \times 2$ matrix $XV_2$, where $V_2$ contains the first 2 columns of $V$ (i.e., project each observation onto the first two principal component directions).

Which statement is most TRUE?

\begin{enumerate}[(A)]
\item Both approaches give identical 2D representations since $U\Sigma = XV$.
\item Approach 1 gives the optimal rank-2 approximation to $X$ in Frobenius norm, while Approach 2 gives the 2D representation that maximizes captured variance, but these are different things.
\item The approaches yield different matrices: $U_2\Sigma_2 \neq XV_2$ because one uses left singular vectors and the other uses right singular vectors.
\item Approach 1 produces vectors in $\mathbb{R}^n$ while Approach 2 produces vectors in $\mathbb{R}^d$, so they cannot be compared.
\item Both approaches produce the same $n \times 2$ matrix of 2D observation coordinates, and both optimally capture the variance structure in a 2D subspace.
\end{enumerate}

% \textbf{Correct Answer:} E
% \textbf{Explanation:} This tests whether students understand the equivalence of different formulations:
% - $U\Sigma V^T$ is the SVD, so $X = U\Sigma V^T$
% - Therefore $XV = U\Sigma V^T V = U\Sigma$ (since $V^TV = I$)
% - Keeping first 2 columns: $U_2\Sigma_2 = XV_2$, where subscript 2 means first 2 columns
% Both approaches yield the SAME $n \times 2$ matrix. This matrix:
% 1) Gives the best rank-2 approximation: $X_2 = U_2\Sigma_2 V_2^T$ minimizes $\|X - X_2\|_F$ (Eckart-Young from Week 10)
% 2) Maximizes variance captured: the columns are projections onto principal directions that explain the most variance (from Week 11)
% 3) Projects observations optimally onto a 2D subspace
% KEY SYNTHESIS: SVD's optimal low-rank approximation and PCA's variance maximization are THE SAME THING for centered data. The geometric picture (ellipsoid with largest semi-axes) corresponds to the statistical picture (directions of maximum variance).
% \textbf{Partial credit:} A (recognizes equivalence but explanation is incomplete)
% \textbf{Partial credit:} B (recognizes both are optimal but claims they're different things, missing that they're the same)
% \textbf{Wrong:} C (both yield identical $n \times 2$ matrices; the identity $XV_2 = U_2\Sigma_2$ follows from $V^TV = I$)
% \textbf{Wrong:} D (both produce $n \times 2$ matrices; confusion about dimensions)
% \textbf{Note:} This is a deep synthesis question showing SVD and PCA are two perspectives on the same mathematical structure

\divider

%%%%%%%%%%%%%%%%%%%%%%%%%%%%%%%%%%%%%%%%%%%%%%%%%%%%%%%%%%%%%%%%%
\newpage
%%%%%%%%%%%%%%%%%%%%%%%%%%%%%%%%%%%%%%%%%%%%%%%%%%%%%%%%%%%%%%%%%

{\bf PROBLEM 4:}
% WEEK = 10
% TOPICS = Singular values ordering, operations on matrices, basic SVD properties

A matrix $A \in \mathbb{R}^{m \times n}$ has singular values $\sigma_1 = 5$, $\sigma_2 = 3$, $\sigma_3 = 1$, and $\sigma_k = 0$ for $k>3$. 

Which of the following statements is most TRUE?

\begin{enumerate}[(A)]
\item Matrix $2A$ has singular values 10, 6, 2, with all others zero, but they should be reordered as 2, 6, 10 to maintain the SVD convention.
\item Matrix $A^T$ has singular values $\sigma_1=1, \sigma_2=3, \sigma_3=5$, with all others zero.
\item Matrix $A^TA$ has eigenvalues 25, 9, 1, with the remaining eigenvalues zero.
\item The largest singular value of $AA^T$ is 25.
\item Matrices $A^TA$ and $AA^T$ have different nonzero eigenvalues.
\end{enumerate}

% \textbf{Correct Answer:} C
% \textbf{Explanation:} Singular values are ALWAYS arranged in descending order: $\sigma_1 \geq \sigma_2 \geq \cdots \geq 0$. This is a convention, not a consequence of any operation. The key relationships: (1) $A$ and $A^T$ have identical singular values (same set, same order). (2) If $A$ has singular values $\{\sigma_i\}$, then $cA$ has singular values $\{|c|\sigma_i\}$ in the same descending order. (3) The eigenvalues of $A^TA$ are $\sigma_i^2$, so $25, 9, 1, 0, 0, \ldots$ in descending order. (4) Both $A^TA$ and $AA^T$ share the same NONZERO eigenvalues (namely $\sigma_i^2$ for $i \leq r = \rank(A)$), though they may have different numbers of zero eigenvalues depending on dimensions.
% \textbf{Partial credit:} None
% \textbf{Trick answer:} A (singular values of $2A$ are correctly computed as $10, 6, 2$, but they do NOT need reordering—they're already in descending order by construction)
% \textbf{Trick answer:} B ($A^T$ has the SAME singular values as $A$, in the same descending order; transpose doesn't reverse the ordering)
% \textbf{Trick answer:} D ($AA^T$ is an $m \times m$ matrix; its eigenvalues are $\sigma_i^2 = 25, 9, 1, 0, \ldots$, so its largest eigenvalue is 25, but this is an eigenvalue, not a singular value of $AA^T$. The singular values of $AA^T$ would be the square roots of eigenvalues of $(AA^T)^T(AA^T) = AA^TAA^T$)
% \textbf{Trick answer:} E (classic misconception: $A^TA$ and $AA^T$ always share the same nonzero eigenvalues, which are $\sigma_1^2, \ldots, \sigma_r^2$)

\divider

{\bf PROBLEM 5:}
% WEEK = 10
% TOPICS = Pseudoinverse construction via SVD, Σ^† formula

A matrix $A \in \mathbb{R}^{4 \times 3}$ has singular value decomposition $A = U\Sigma V^T$ where $U \in \mathbb{R}^{4 \times 4}$, $V \in \mathbb{R}^{3 \times 3}$ are orthogonal, and:
$$\Sigma = \begin{bmatrix} 5 & 0 & 0 \\ 0 & 2 & 0 \\ 0 & 0 & 0 \\ 0 & 0 & 0 \end{bmatrix} \in \mathbb{R}^{4 \times 3}$$

What is the pseudoinverse $A^\dagger$ of this matrix?

\begin{enumerate}[(A)]
\item $A^\dagger = U\Sigma^{-1} V^T$ where $\Sigma^{-1} = \begin{bmatrix} 1/5 & 0 & 0 \\ 0 & 1/2 & 0 \\ 0 & 0 & \text{undefined} \\ 0 & 0 & \text{undefined} \end{bmatrix}$
\item $A^\dagger = V\Sigma^T U^T$ where $\Sigma^T = \begin{bmatrix} 5 & 0 & 0 & 0 \\ 0 & 2 & 0 & 0 \\ 0 & 0 & 0 & 0 \end{bmatrix}$
\item $A^\dagger = V\Sigma^\dagger U^T$ where $\Sigma^\dagger = \begin{bmatrix} 1/5 & 0 & 0 & 0 \\ 0 & 1/2 & 0 & 0 \\ 0 & 0 & 0 & 0 \end{bmatrix}$
\item $A^\dagger = U^T\Sigma^\dagger V$ where $\Sigma^\dagger = \begin{bmatrix} 1/5 & 0 & 0 & 0 \\ 0 & 1/2 & 0 & 0 \\ 0 & 0 & 0 & 0 \end{bmatrix}$
\item $A^\dagger = V\Sigma^\dagger U^T$ where $\Sigma^\dagger = \begin{bmatrix} 1/5 & 0 & 0 \\ 0 & 1/2 & 0 \\ 0 & 0 & 0 \\ 0 & 0 & 0 \end{bmatrix}$
\end{enumerate}

% \textbf{Correct Answer:} C
% \textbf{Explanation:} The pseudoinverse formula is $A^\dagger = V\Sigma^\dagger U^T$. The pseudoinverse of the diagonal matrix $\Sigma$ is constructed by: (1) transposing $\Sigma$ to get a $3 \times 4$ matrix, then (2) replacing each nonzero diagonal entry $\sigma_i$ with $1/\sigma_i$, and (3) leaving zero entries as zero. Thus:
% $$\Sigma^\dagger = \begin{bmatrix} 1/5 & 0 & 0 & 0 \\ 0 & 1/2 & 0 & 0 \\ 0 & 0 & 0 & 0 \end{bmatrix} \in \mathbb{R}^{3 \times 4}$$
% Note the dimensions: $A$ is $4 \times 3$, so $A^\dagger$ must be $3 \times 4$. We have $V$ is $3 \times 3$, $\Sigma^\dagger$ is $3 \times 4$, and $U^T$ is $4 \times 4$, giving $(3 \times 3)(3 \times 4)(4 \times 4) = 3 \times 4$. ✓ The key insight: we don't try to "invert" $\Sigma$ (which is rectangular and rank-deficient), but rather construct its pseudoinverse by transposing and inverting only the nonzero entries.
% \textbf{Partial credit:} None
% \textbf{Trick answer:} A (trying to directly invert $\Sigma$ doesn't work because $\Sigma$ is rectangular and has zero entries; the pseudoinverse handles this by transposing first)
% \textbf{Trick answer:} B (just transposing $\Sigma$ doesn't give the pseudoinverse; you must also invert the nonzero entries)
% \textbf{Trick answer:} D (incorrect ordering; the pseudoinverse is $V\Sigma^\dagger U^T$, not $U^T\Sigma^\dagger V$, which comes from reversing $A = U\Sigma V^T$)
% \textbf{Trick answer:} E (zero singular values remain zero in $\Sigma^\dagger$; they are NOT replaced with 1 or any other value)

\divider


{\bf PROBLEM 6:}
% WEEK = 10
% TOPICS = Fundamental subspaces via SVD, left vs right singular vectors, geometric interpretation

A data matrix $A \in \mathbb{R}^{100 \times 50}$ represents measurements where rows are observations and columns are features. The SVD $A = U\Sigma V^T$ reveals that only the first 10 singular values are nonzero: $\sigma_1 \geq \sigma_2 \geq \cdots \geq \sigma_{10} > 0$ and $\sigma_{11} = \cdots = \sigma_{50} = 0$.

An analyst wants to understand the fundamental subspaces. Which statement is most TRUE?

\begin{enumerate}[(A)]
\item The first 10 right singular vectors $\mathbf{v}_1, \ldots, \mathbf{v}_{10}$ span the column space of $A$, capturing the 10 most important patterns in the observations.
\item The first 10 left singular vectors $\mathbf{u}_1, \ldots, \mathbf{u}_{10}$ span a 10-dimensional subspace of feature space where all data variation occurs.
\item The right singular vectors $\mathbf{v}_{11}, \ldots, \mathbf{v}_{50}$ form a basis for $\ker(A)$, representing 40 redundant feature combinations.
\item The left singular vectors $\mathbf{u}_{11}, \ldots, \mathbf{u}_{100}$ span $\ker(A^T)$, representing observation directions that have no correlation with any features.
\item Both (C) and (D) are correct.
\end{enumerate}

% \textbf{Correct Answer:} E
% \textbf{Explanation:} This tests understanding of which singular vectors span which fundamental subspaces. For $A \in \mathbb{R}^{100 \times 50}$ with rank 10:
% - Right singular vectors $\mathbf{v}_i \in \mathbb{R}^{50}$ (feature space): $\{\mathbf{v}_1, \ldots, \mathbf{v}_{10}\}$ span $\text{im}(A^T)$ (NOT $\text{im}(A)$), and $\{\mathbf{v}_{11}, \ldots, \mathbf{v}_{50}\}$ span $\ker(A)$. Statement C is correct.
% - Left singular vectors $\mathbf{u}_i \in \mathbb{R}^{100}$ (observation space): $\{\mathbf{u}_1, \ldots, \mathbf{u}_{10}\}$ span $\text{im}(A)$, and $\{\mathbf{u}_{11}, \ldots, \mathbf{u}_{100}\}$ span $\ker(A^T)$. Statement D is correct.
% The kernel $\ker(A)$ has dimension $50 - 10 = 40$, representing linear combinations of features that always give zero when $A$ acts on them (redundant feature directions). The left null space $\ker(A^T)$ has dimension $100 - 10 = 90$, representing observation directions orthogonal to the column space.
% \textbf{Partial credit:} C (correctly identifies the kernel basis but misses that (D) is also correct)
% \textbf{Partial credit:} D (correctly identifies the left null space but misses that (C) is also correct)
% \textbf{Trick answer:} A (confuses $\text{im}(A^T)$ with $\text{im}(A)$; right singular vectors span $\text{im}(A^T)$ in the domain, not $\text{im}(A)$ in the codomain)
% \textbf{Trick answer:} B (left singular vectors live in observation space $\mathbb{R}^{100}$, not feature space $\mathbb{R}^{50}$; they span the column space $\text{im}(A)$, which is 10-dimensional)

\divider

{\bf PROBLEM 7:}
% WEEK = 10
% TOPICS = Matrix dimensions in SVD, compatibility, parameter counting

A machine learning engineer wants to store a compressed representation of matrix $A \in \mathbb{R}^{500 \times 300}$ using its rank-$k$ SVD approximation $A_k = \sum_{i=1}^k \sigma_i \mathbf{u}_i\mathbf{v}_i^T$.

To store $A_k$, they need to save:
\begin{itemize}
\item $k$ singular values $\{\sigma_1, \ldots, \sigma_k\}$
\item $k$ left singular vectors $\{\mathbf{u}_1, \ldots, \mathbf{u}_k\}$
\item $k$ right singular vectors $\{\mathbf{v}_1, \ldots, \mathbf{v}_k\}$
\end{itemize}

How many total scalar values must be stored, and for which value of $k$ does the storage equal the original matrix $A$?

\begin{enumerate}[(A)]
\item $k(500 + 300 + 1)$ values; storage equals original when $k = 500$
\item $k(500 + 300 + 1)$ values; storage equals original when $k = 300$
\item $k(500 + 300)$ values; storage equals original when $k = 187$ (approximately)
\item $k(500 \cdot 300)$ values; storage never equals original since SVD is always compressed
\item $k(500 + 300 + 1)$ values; storage equals original when $k = 187$ (approximately)
\end{enumerate}

% \textbf{Correct Answer:} E
% \textbf{Explanation:} For rank-$k$ approximation of $A \in \mathbb{R}^{500 \times 300}$:
% - $k$ singular values (each is 1 scalar): $k$ values
% - $k$ left singular vectors $\mathbf{u}_i \in \mathbb{R}^{500}$: $500k$ values
% - $k$ right singular vectors $\mathbf{v}_i \in \mathbb{R}^{300}$: $300k$ values
% Total storage: $k + 500k + 300k = k(500 + 300 + 1) = 801k$ values.
%
% The original matrix $A$ stores $500 \times 300 = 150000$ values. Storage equals original when:
% $801k = 150000$, so $k = 150000/801 \approx 187.3$
%
% This shows the crossover point: for $k < 187$, SVD compression saves storage; for $k > 187$, it uses more. The maximum useful rank is $\min(500, 300) = 300$, but full rank ($k=300$) would store $801 \times 300 = 240300$ values, which exceeds the original! This reveals why SVD compression is most valuable when the effective rank is much smaller than both dimensions.
% \textbf{Partial credit:} None
% \textbf{Trick answer:} A (correctly counts storage but incorrectly thinks full storage occurs at $k = m = 500$; the maximum rank is $\min(m,n) = 300$)
% \textbf{Trick answer:} B (correctly counts storage but thinks full storage = max rank; actually at max rank, SVD stores MORE than original due to separate storage of vectors)
% \textbf{Trick answer:} C (forgets to count the $k$ singular values themselves; singular values are scalars that must also be stored)
% \textbf{Trick answer:} D (incorrectly thinks each component requires full $m \times n$ storage; the outer product $\mathbf{u}_i\mathbf{v}_i^T$ need not be stored explicitly)

\divider

%%%%%%%%%%%%%%%%%%%%%%%%%%%%%%%%%%%%%%%%%%%%%%%%%%%%%%%%%%%%%%%%%
\newpage
%%%%%%%%%%%%%%%%%%%%%%%%%%%%%%%%%%%%%%%%%%%%%%%%%%%%%%%%%%%%%%%%%

{\bf PROBLEM 8:}
% WEEK = 10
% TOPICS = SVD of orthogonal matrices, singular values, geometric interpretation

What can be said about the SVD of an orthogonal matrix $Q \in \mathbb{R}^{n \times n}$?. 

\begin{enumerate}[(A)]
\item All singular values equal 1, so $Q = U\Sigma V^T$ where $\Sigma = I$, meaning $Q = UV^T$ is a product of two orthogonal matrices.
\item The singular values depend on the specific matrix $Q$; for example, a rotation matrix has different singular values than a reflection matrix.
\item All singular values equal 1, and because of this the SVD is not unique.
\item Orthogonal matrices do not have a uniquely defined SVD because they preserve lengths, which contradicts the stretching interpretation of singular values.
\item The largest singular value is 1, but the other singular values can be less than 1 depending on how much $Q$ ``compresses'' certain directions.
\end{enumerate}

% \textbf{Correct Answer:} A
% \textbf{Explanation:} For any orthogonal matrix $Q$, we have $Q^T Q = I$. The singular values of $Q$ are defined as $\sigma_i = \sqrt{\lambda_i}$ where $\lambda_i$ are eigenvalues of $Q^T Q = I$. Since the eigenvalues of the identity matrix are all 1, we get $\sigma_i = \sqrt{1} = 1$ for all $i$. Therefore, in the SVD $Q = U\Sigma V^T$, we have $\Sigma = I$, giving $Q = U I V^T = UV^T$. This shows that every orthogonal matrix can be written as a product of two orthogonal matrices. Geometrically, orthogonal matrices preserve lengths (no stretching or compression), which is exactly what having all singular values equal to 1 means. The SVD decomposition $Q = UV^T$ represents $Q$ as a composition of two orthogonal transformations.
% \textbf{Partial credit:} None
% \textbf{Trick answer:} B (doesn't recognize that the defining property $Q^T Q = I$ immediately forces all singular values to be 1)
% \textbf{Trick answer:} C (correctly identifies that all singular values are 1, but incorrectly concludes non-uniqueness
% \textbf{Trick answer:} D (misunderstands the SVD; preserving lengths means singular values are 1, not that SVD doesn't exist)
% \textbf{Trick answer:} E (backwards reasoning: orthogonal matrices preserve ALL lengths, so no direction is compressed; all singular values must be 1)

\divider


{\bf PROBLEM 9:}
% WEEK = 11
% TOPICS = Covariance vs Correlation PCA, preprocessing and scaling

An engineer monitors a chemical reactor with three sensors measuring temperature ($T$, in °C), pressure ($P$, in kPa), and flow rate ($F$, in mL/min). A week of measurements yields:
%
$$\text{Temperature: } T \in [180, 220] \text{ °C}$$
$$\text{Pressure: } P \in [2000, 2500] \text{ kPa}$$
$$\text{Flow rate: } F \in [5, 15] \text{ mL/min}$$

The engineer performs covariance PCA on the centered data and finds that the first principal component is approximately:
$$\mathbf{v}_1 \approx (0.002, 0.998, 0.005)^T$$

Which statement best explains this result?

\begin{enumerate}[(A)]
\item Pressure dominates the first principal component because it has the largest absolute variance, masking potentially important patterns in temperature and flow.
\item The principal component correctly identifies pressure as the most important variable since it has the widest measurement range.
\item This result suggests the three variables are uncorrelated, with pressure varying independently of temperature and flow.
\item The small coefficients for temperature and flow indicate these variables contribute negligible information and can be discarded.
\item Covariance PCA automatically accounts for different measurement scales, so this result reveals genuine physical dominance of pressure variation.
\end{enumerate}

% \textbf{Correct Answer:} A
% \textbf{Explanation:} Covariance PCA is sensitive to the scale of variables. Pressure varies on the order of hundreds of kPa, while temperature varies by tens of degrees and flow by ~10 mL/min. The sample variance of pressure will be much larger (proportional to the square of the range), causing the first PC to align almost entirely with the pressure direction. This doesn't mean pressure is genuinely "more important" - it's an artifact of the different measurement units. To reveal patterns independent of scale, the engineer should use correlation PCA (standardize variables to unit variance before PCA) or carefully rescale the variables.
% \textbf{Partial credit:} None
% \textbf{Wrong:} B (confuses large range with importance; this is precisely the problem with covariance PCA when scales differ)
% \textbf{Wrong:} C (the PC structure doesn't directly reveal correlation; variables could be highly correlated but the large-scale variable still dominates)
% \textbf{Wrong:} D (the small coefficients are due to scaling issues, not lack of information)
% \textbf{Trick answer:} E (contradicts the fundamental issue with covariance PCA - it does NOT automatically account for scale)

\divider

%%%%%%%%%%%%%%%%%%%%%%%%%%%%%%%%%%%%%%%%%%%%%%%%%%%%%%%%%%%%%%%%%
\newpage
%%%%%%%%%%%%%%%%%%%%%%%%%%%%%%%%%%%%%%%%%%%%%%%%%%%%%%%%%%%%%%%%%


{\bf PROBLEM 10:}
% WEEK = 11
% TOPICS = Interpreting principal components, physical meaning

Four vibration sensors measure vertical displacement (in mm) on a suspension bridge at locations: North anchor, North midspan, South midspan, South anchor. After centering the data, (covariance) PCA reveals the following eigenvectors/eigenvalues of $X^TX$:

$$\mathbf{v}_1 = \frac{1}{2}\begin{pmatrix} 1 \\ 1 \\ 1 \\ 1 \end{pmatrix}, \quad \mathbf{v}_2 = \frac{1}{\sqrt{2}}\begin{pmatrix} 1 \\ 0 \\ 0 \\ -1 \end{pmatrix}, \quad \mathbf{v}_3 = \frac{1}{\sqrt{2}}\begin{pmatrix} 0 \\ 1 \\ -1 \\ 0 \end{pmatrix}$$

with eigenvalues $\lambda_1 = 36$, $\lambda_2 = 9$, $\lambda_3 = 4$, $\lambda_4 = 1$ (in mm²).

Which interpretation is most consistent with these results?

\begin{enumerate}[(A)]
\item $\mathbf{v}_1$ represents uniform vertical motion of the entire bridge (all sensors move together), capturing 72\% of total variance.
\item $\mathbf{v}_2$ represents torsional twisting where the North side moves opposite to the South side, while midspan points remain stationary.
\item $\mathbf{v}_3$ represents antisymmetric bending where the two midspan points move in opposite directions while anchors remain fixed.
\item All of the above (A, B, and C) are correct interpretations of the principal components.
\item None of the above.
\end{enumerate}

% \textbf{Correct Answer:} D
% \textbf{Explanation:} Each principal component reveals a distinct vibration mode:
% - $\mathbf{v}_1 = (1,1,1,1)/2$: All four sensors move together with equal weight → uniform vertical displacement. Proportion: $36/50 = 72\%$.
% - $\mathbf{v}_2 = (1,0,0,-1)/\sqrt{2}$: North anchor and South anchor move in opposite directions (+1 vs -1), midspan sensors have zero weight → end-to-end "tilting" or rotational mode.
% - $\mathbf{v}_3 = (0,1,-1,0)/\sqrt{2}$: The two midspan sensors move in opposition while anchors have zero weight → midspan antisymmetric bending.
% These interpretations match physical vibration modes of bridges. The large $\lambda_1$ indicates most variation is uniform heaving.
% \textbf{Partial credit:} A, B, or C individually (each is correct but incomplete)
% \textbf{Wrong:} E (PCs often have clear physical interpretations, especially in spatial/temporal data)
% \textbf{Note:} This tests whether students understand that PC vectors show HOW variables combine, with entries representing weights/contributions

\divider 



{\bf PROBLEM 11:}
% WEEK = 10
% TOPICS = SVD revealing Fundamental Theorem structure, isomorphism between subspaces

The SVD $A = U\Sigma V^T$ reveals that $A$ acts as an isomorphism (scaled by singular values) between certain fundamental subspaces. For matrix $A \in \mathbb{R}^{m \times n}$ with rank $r$, which statement is most TRUE?

\begin{enumerate}[(A)]
\item $A$ maps $\ker(A)$ isomorphically onto $\ker(A^T)$ by the relationship $A\mathbf{v}_i = \sigma_i \mathbf{u}_i$ for $i > r$.
\item $A$ maps $\text{im}(A^T)$ isomorphically onto $\text{im}(A)$, with the isomorphism given by scaling: $A\mathbf{v}_i = \sigma_i \mathbf{u}_i$ for $i \leq r$.
\item $A$ maps $\mathbb{R}^n$ isomorphically onto $\mathbb{R}^m$ whenever $m = n$, regardless of rank.
\item $A$ maps both $\ker(A)$ and $\text{im}(A^T)$ isomorphically onto $\text{im}(A)$, explaining why $\dim(\ker(A)) + \dim(\text{im}(A^T)) = \dim(\text{im}(A))$.
\item The restriction of $A$ to $\ker(A)$ is an isomorphism onto $\ker(A^T)$ because both have dimension $n - r$.
\end{enumerate}

% \textbf{Correct Answer:} B
% \textbf{Explanation:} This tests the deep geometric insight from the Fundamental Theorem via SVD. The key observation: $A$ collapses $\ker(A)$ to zero but acts as an isomorphism on $(\ker A)^\perp = \text{im}(A^T)$. Specifically:
% - For $i \leq r$ (nonzero singular values): $A\mathbf{v}_i = \sigma_i \mathbf{u}_i$. The vectors $\{\mathbf{v}_1, \ldots, \mathbf{v}_r\}$ form an orthonormal basis for $\text{im}(A^T)$, and $\{\mathbf{u}_1, \ldots, \mathbf{u}_r\}$ form an orthonormal basis for $\text{im}(A)$.
% - $A$ restricted to $\text{im}(A^T)$ gives a bijection onto $\text{im}(A)$: it's injective (kernel restricted to $\text{im}(A^T)$ is trivial) and surjective (by definition of image). The singular values $\sigma_i$ provide the scaling factors for this isomorphism.
% - For $i > r$ (zero singular values): $A\mathbf{v}_i = \mathbf{0}$, so $A$ collapses $\ker(A) = \text{span}\{\mathbf{v}_{r+1}, \ldots, \mathbf{v}_n\}$ to zero.
% This reveals the geometric content of the Fundamental Theorem: $A$ acts as $(\text{im}(A^T) \xrightarrow{\sim} \text{im}(A)) \oplus (\ker(A) \to \{\mathbf{0}\})$.
% \textbf{Partial credit:} None
% \textbf{Trick answer:} A ($A$ maps $\ker(A)$ to zero, not to $\ker(A^T)$; these are unrelated subspaces in different vector spaces)
% \textbf{Trick answer:} C (false even when $m = n$; if rank$(A) < n$, then $A$ is not an isomorphism because it has a nontrivial kernel)
% \textbf{Trick answer:} D (completely wrong dimension formula; rank-nullity says $\dim(\ker(A)) + \dim(\text{im}(A^T)) = n$, not $\dim(\text{im}(A))$)
% \textbf{Trick answer:} E ($A$ restricted to $\ker(A)$ gives the zero map, not an isomorphism; $A$ annihilates $\ker(A)$)

\divider

%%%%%%%%%%%%%%%%%%%%%%%%%%%%%%%%%%%%%%%%%%%%%%%%%%%%%%%%%%%%%%%%%
\newpage
%%%%%%%%%%%%%%%%%%%%%%%%%%%%%%%%%%%%%%%%%%%%%%%%%%%%%%%%%%%%%%%%%


%%%%%%%%%%%%%%%%%%%%%%%%%%%%%%%%%%%%%%%%%%%%%%%%%%%%%%%%%%%%%%%%%
\newpage
%%%%%%%%%%%%%%%%%%%%%%%%%%%%%%%%%%%%%%%%%%%%%%%%%%%%%%%%%%%%%%%%%

{\bf PROBLEM 12:}
% WEEK = 11
% TOPICS = Covariance vs correlation PCA, scree plots, intrinsic dimensionality, scale artifacts

A data scientist analyzes measurements from six sensors monitoring a thermodynamic process. The sensors measure various properties of the system: pressure, temperature, volume, entropy, enthalpy, and work.
The figures below plot (vertically) cumulative explained variance for covariance PCA (left) and correlation PCA (right), versus (horizontally) number of singular values (from 0 to 6). Based on these plots, which interpretation is most appropriate?


% [LEFT FIGURE - Covariance PCA: Cumulative variance plot]
% Points at: (1, 94%), (2, 97%), (3, 98.5%), (4, 99.2%), (5, 99.7%), (6, 100%)
% Nearly flat after first component - looks like 1D data

% [RIGHT FIGURE - Correlation PCA: Cumulative variance plot]  
% Points at: (1, 45%), (2, 75%), (3, 90%), (4, 96%), (5, 99%), (6, 100%)
% Clear elbow at k=3

\begin{center}
\begin{figure}[h]
\includegraphics[width=7in]{2030 Q5 SCREE COV-COR.jpg}
%\caption{[left] for PROBLEM 14 ; [right] for PROBLEM 15.}
\end{figure}
\end{center}
\begin{enumerate}[(A)]
\item Covariance PCA reveals that the process is essentially one-dimensional: a single latent variable explains 94\% of all variation, making $k=1$ the optimal choice for dimensionality reduction.

\item Correlation PCA suggests the process has intrinsic dimension approximately 3, with scale differences among sensors masking this structure in the covariance analysis.

\item Both analyses agree on low intrinsic dimensionality; correlation PCA is always preferable as it preserves the natural units of measurement.

\item The correlation PCA result is misleading because standardizing variables inflates the apparent contribution of low-variance sensors like volume.

\item A rank-2 approximation captures over 75\% of variance in both analyses, so either method supports reducing to 2 dimensions.
\end{enumerate}

% \textbf{Correct Answer:} B
% \textbf{Explanation:} The 94% variance captured by the first covariance PC is a classic symptom of 
% scale mismatch: pressure (measured in thousands of kPa) dominates the variance simply due to units.
% This doesn't indicate true 1-dimensionality—it indicates that one variable has numerically larger 
% values than the others.
%
% Correlation PCA standardizes all variables to unit variance, removing scale artifacts. The resulting
% plot shows a clear elbow at $k=3$: three components capture 90% of standardized variance, suggesting
% the underlying process has approximately 3 degrees of freedom. This is the genuine intrinsic 
% dimensionality, revealed only after accounting for measurement scale.
%
% The covariance result would lead to discarding potentially important process information that 
% happens to be measured in small-magnitude units (like pH ranging 6-8 or dissolved O₂ in mg/L).
%
% \textbf{Partial credit:} None
% \textbf{Trick answer:} A (falls into the scale-artifact trap; high variance ≠ physical importance)
% \textbf{Trick answer:} C (covariance PCA does NOT reveal low dimensionality—it reveals scale dominance; 
%   also, "natural units" is not a valid reason when units are incommensurable)
% \textbf{Trick answer:} D (standardization doesn't "inflate" low-variance variables unfairly—it gives 
%   each variable equal opportunity to contribute, which is appropriate when scales are arbitrary)
% \textbf{Trick answer:} E (technically true that correlation PCA hits 75% at $k=2$, but covariance PCA 
%   hits 97% at $k=2$, so "both analyses" is misleading; more importantly, this misses the key insight 
%   about WHICH analysis to trust)

\divider


%%%%%%%%%%%%%%%%%%%%%%%%%%%%%%%%%%%%%%%%%%%%%%%%%%%%%%%%%%%%%%%%%
\newpage
%%%%%%%%%%%%%%%%%%%%%%%%%%%%%%%%%%%%%%%%%%%%%%%%%%%%%%%%%%%%%%%%%

{\bf PROBLEM 13:}
% WEEK = 11
% TOPICS = Choosing number of components, singular value decay patterns

Two centered data matrices are analyzed via covariance PCA:

\textbf{Dataset A:} $X_A \in \mathbb{R}^{200 \times 10}$ has all nonzero singular values:
$$\sigma_1 = 24, \quad \sigma_2 = 22, \quad \sigma_3 = 20, \quad \sigma_4 = 18, \quad \sigma_5 = 16, \ldots$$

\textbf{Dataset B:} $X_B \in \mathbb{R}^{200 \times 10}$ has all nonzero singular values:
$$\sigma_1 = 50, \quad \sigma_2 = 25, \quad \sigma_3 = 8, \quad \sigma_4 = 3, \quad \sigma_5 = 1, \ldots$$

Which statement best characterizes the difference between these datasets?

\begin{enumerate}[(A)]
\item Dataset A has clearer low-dimensional structure.
\item Dataset B has clearer low-dimensional structure.
\item Dataset A is better for dimensionality reduction based on better total variance.
\item Both datasets have similar intrinsic dimensionality since all singular values are positive.
\item Dataset B should have been handled via correlation PCA.
\end{enumerate}

% \textbf{Correct Answer:} B
% \textbf{Explanation:} The key insight is interpreting singular value decay. Dataset A shows gradual, nearly uniform decay (24, 22, 20, 18, ...) suggesting variation is spread across many dimensions with no clear dominant low-dimensional structure. Dataset B shows rapid decay (50, 25, 8, 3, 1, ...) with a sharp drop after the first 2-3 values, indicating the data lies approximately in a low-dimensional subspace. For PCA/dimensionality reduction, rapid decay is desirable - it means a few components capture most variation. This is the "elbow" pattern: large early values, then a sharp drop to small values. Dataset B would need fewer components to achieve any given variance threshold.
% \textbf{Partial credit:} None
% \textbf{Wrong:} A (uniform singular values mean NO clear low-dimensional structure - variation is spread equally)
% \textbf{Wrong:} C (total variance doesn't determine reducibility; the distribution/decay pattern matters)
% \textbf{Wrong:} D (intrinsic dimensionality depends on data structure, not matrix dimensions)
% \textbf{Trick answer:} E (rapid decay means FEWER components needed, not more)

\divider


{\bf PROBLEM 14:}
% WEEK = 10
% TOPICS = Four fundamental subspaces via SVD, kernel, image, orthogonal complements

A matrix $A \in \mathbb{R}^{7 \times 5}$ has rank 3 with SVD $A = U\Sigma V^T$.

Which approach yields a basis for the coimage $\text{coim}(A) = (\ker A)^\perp$?

\begin{enumerate}[(A)]
\item Right singular vectors $\{\mathbf{v}_1, \mathbf{v}_2, \mathbf{v}_3, \mathbf{v}_4, \mathbf{v}_5\}$
\item Left singular vectors $\{\mathbf{u}_1, \mathbf{u}_2, \mathbf{u}_3\}$
\item Right singular vectors $\{\mathbf{v}_1, \mathbf{v}_2, \mathbf{v}_3\}$
\item Left singular vectors $\{\mathbf{u}_4, \mathbf{u}_5, \mathbf{u}_6, \mathbf{u}_7\}$
\item Right singular vectors $\{\mathbf{v}_4, \mathbf{v}_5\}$
\end{enumerate}

% \textbf{Correct Answer:} C
% \textbf{Explanation:} The coimage $\text{coim}(A) = (\ker A)^\perp$ is the orthogonal complement of the kernel in the domain $\mathbb{R}^5$. By the Fundamental Theorem, $(\ker A)^\perp = \text{im}(A^T)$. The SVD reveals that $\text{im}(A^T)$ is spanned by right singular vectors corresponding to nonzero singular values. Since rank$(A) = 3$, we have exactly 3 nonzero singular values, so $\text{coim}(A) = \text{span}\{\mathbf{v}_1, \mathbf{v}_2, \mathbf{v}_3\}$.
% \textbf{Partial credit:} None
% \textbf{Trick answer:} B (left singular vectors live in codomain $\mathbb{R}^7$, not domain $\mathbb{R}^5$)
% \textbf{Trick answer:} D (wrong space and wrong singular values)
% \textbf{Trick answer:} E (these span $\ker(A)$, not its orthogonal complement)
% \textbf{Trick answer:} A (includes both $\ker(A)$ and $(\ker A)^\perp$, giving all of $\mathbb{R}^5$ rather than just the coimage)

\divider

{\bf PROBLEM 15:}
% WEEK = 10
% TOPICS = Eigenvalues vs singular values, symmetric matrices, general square matrices

A $3 \times 3$ matrix $A$ has eigenvalues $\lambda_1 = 5$, $\lambda_2 = -3$, $\lambda_3 = 1$.
%
Which statement about the singular values of $A$ is most TRUE?

\begin{enumerate}[(A)]
\item The singular values are $5, 3, 1$.
\item The singular values are $25, 9, 1$.
\item The singular values are $\sqrt{5}, \sqrt{3}, 1$.
\item The singular values cannot be determined from eigenvalues alone.
\item The singular values of a {\em square} matrix equal the eigenvalues for any square matrix.
\end{enumerate}

% \textbf{Correct Answer:} D
% \textbf{Explanation:} Eigenvalues and singular values measure fundamentally different things:
% - Eigenvalues come from $A\mathbf{v} = \lambda\mathbf{v}$ and can be negative or complex
% - Singular values come from $A^TA$ (or the SVD) and are always real and non-negative
% 
% For SYMMETRIC matrices: $A = A^T$ implies $A^TA = A^2$, so the eigenvalues of $A^TA$ are 
% $\lambda_i^2$, giving singular values $\sigma_i = |\lambda_i|$. Thus: $\sigma_1 = 5$, $\sigma_2 = 3$, $\sigma_3 = 1$.
% 
% For NON-SYMMETRIC matrices: the eigenvalues of $A$ and $A^TA$ are unrelated in general.
% A non-symmetric matrix with these eigenvalues could have completely different singular values.
% Example: a shear matrix can have repeated eigenvalue 1 but singular values far from 1.
% 
% \textbf{Partial credit:} A (correct answer IF the matrix is symmetric, but misses the conditional)
% \textbf{Trick answer:} B (confuses eigenvalues of $A^TA$ with singular values; singular values are $\sqrt{\lambda(A^TA)}$, not $\lambda^2$)
% \textbf{Trick answer:} C (inverts the relationship; singular values are $|\lambda_i|$ for symmetric matrices, not $\sqrt{|\lambda_i|}$)
% \textbf{Trick answer:} E (false: eigenvalues can be negative or complex, but singular values are always $\geq 0$; immediately contradicted by $\lambda_2 = -3$)

\divider


%%%%%%%%%%%%%%%%%%%%%%%%%%%%%%%%%%%%%%%%%%%%%%%%%%%%%%%%%%%%%%%%%
\newpage
%%%%%%%%%%%%%%%%%%%%%%%%%%%%%%%%%%%%%%%%%%%%%%%%%%%%%%%%%%%%%%%%%



{\bf PROBLEM 16:}
% WEEK = 11
% TOPICS = Effect of orthogonal transformations, invariance properties, creative application

Consider two analysts working with the same centered data matrix $X \in \mathbb{R}^{100 \times 6}$ that has singular values $\sigma_1 = 20, \sigma_2 = 15, \sigma_3 = 8, \sigma_4 = 5, \sigma_5 = 2, \sigma_6 = 1$.

\textbf{Analyst A} performs PCA directly on $X$.

\textbf{Analyst B} first applies an orthogonal transformation $Q \in \mathbb{R}^{6 \times 6}$ to get $Y = XQ$, then performs PCA on $Y$.

How do their results compare?

\begin{enumerate}[(A)]
\item Analyst B's principal components are related to Analyst A's by the transformation: $\mathbf{w}_k = Q^T\mathbf{v}_k$, where $\mathbf{v}_k$ are A's principal components.
\item Both analysts obtain identical singular values: $\{20, 15, 8, 5, 2, 1\}$, but their principal component directions differ.
\item Analyst A's results are superior because applying $Q$ distorts the covariance structure of the data.
\item The analyses are equivalent only if $Q$ is a rotation (determinant +1), but not if $Q$ includes reflections (determinant -1).
\item Both analyses capture the same total variance and the same variance proportions, but the ordering of principal components may differ.
\end{enumerate}

% \textbf{Correct Answer:} B
% \textbf{Explanation:} Key insight: Orthogonal transformations preserve singular values. If $X = U\Sigma V^T$, then $Y = XQ = U\Sigma V^TQ$. The singular values of $Y$ are found from $Y^TY = Q^TV\Sigma^2 V^TQ$. Since $Q$ is orthogonal, this is similar to $V\Sigma^2 V^T$, preserving eigenvalues. Thus singular values of $Y$ equal those of $X$.
% However, the principal components differ: for $Y$, the PCs are columns of $W$ where $Y = U\Sigma W^T$, giving $W = Q^TV$, so $\mathbf{w}_k = Q^T\mathbf{v}_k$.
% CONCEPTUAL POINT: The singular value spectrum (and thus variance captured) is invariant under orthogonal transformations of the columns. This makes sense: rotating the coordinate system doesn't change the intrinsic dimensionality or variance structure of the data - just the coordinate representation of the principal directions.
% \textbf{Partial credit:} A (correct relationship but doesn't state that singular values are preserved)
% \textbf{Wrong:} C (orthogonal transformations preserve geometric structure; neither is "superior")
% \textbf{Wrong:} D (all orthogonal transformations preserve singular values, regardless of determinant)
% \textbf{Trick answer:} E (the variance proportions are identical, but the ordering remains the same - both get $\sigma_1 = 20$ as largest, etc.)

\divider


{\bf PROBLEM 17:}
% WEEK = 11
% TOPICS = Correlation as cosine similarity, orthogonality, angle interpretation

Four feature vectors $(X_1, X_2, X_3, X_4)$ are extracted from centered sensor data. The covariance and correlation matrices are:

$$[\Sigma] = \begin{bmatrix}
4 & 3 & -4.2 & 6.4 \\
3 & 25 & -6 & -12 \\
-4.2 & -6 & 9 & 1.2 \\
6.4 & -12 & 1.2 & 16
\end{bmatrix}
\quad : \quad
[R] = \begin{bmatrix}
1.0 & 0.3 & -0.7 & 0.8 \\
0.3 & 1.0 & -0.4 & -0.6 \\
-0.7 & -0.4 & 1.0 & 0.1 \\
0.8 & -0.6 & 0.1 & 1.0
\end{bmatrix}$$

Which pairs of feature vectors have directions within $30^\circ$ of being orthogonal?

\begin{enumerate}[(A)]
\item $(X_1, X_3)$ and $(X_2, X_4)$
\item $(X_1, X_2)$ and $(X_3, X_4)$
\item $(X_2, X_3)$ only
\item $(X_1, X_3)$, $(X_2, X_3)$, and $(X_2, X_4)$
\item $(X_1, X_2)$, $(X_2, X_3)$, and $(X_3, X_4)$
\end{enumerate}

% \textbf{Correct Answer:} E
% \textbf{Explanation:} "Within 30° of orthogonal" means the angle θ satisfies $60° < \theta < 120°$, which requires:
% $$\cos(60°) < \cos(\theta) < \cos(120°) \quad \Rightarrow \quad -0.5 < \rho < 0.5$$
% Since ρ = cos(θ) for centered vectors, we need $|\rho| < 0.5$. Checking each pair:
% - $(X_1, X_2)$: $\rho_{12} = 0.3$, so $|\rho| = 0.3 < 0.5$ → angle ≈ 73° ✓
% - $(X_1, X_3)$: $\rho_{13} = -0.7$, so $|\rho| = 0.7 > 0.5$ → angle ≈ 134° ✗
% - $(X_1, X_4)$: $\rho_{14} = 0.8$, so $|\rho| = 0.8 > 0.5$ → angle ≈ 37° ✗
% - $(X_2, X_3)$: $\rho_{23} = -0.4$, so $|\rho| = 0.4 < 0.5$ → angle ≈ 114° ✓
% - $(X_2, X_4)$: $\rho_{24} = -0.6$, so $|\rho| = 0.6 > 0.5$ → angle ≈ 127° ✗
% - $(X_3, X_4)$: $\rho_{34} = 0.1$, so $|\rho| = 0.1 < 0.5$ → angle ≈ 84° ✓
% Pairs within 30° of orthogonal: $(X_1, X_2)$, $(X_2, X_3)$, $(X_3, X_4)$.
% KEY DISTRACTION: Negative correlation doesn't automatically mean near-orthogonal. $(X_1, X_3)$ has $\rho = -0.7$ (angle 134°) and $(X_2, X_4)$ has $\rho = -0.6$ (angle 127°), both too far from 90°. Similarly, small covariance doesn't guarantee near-orthogonality - must check correlation.
% \textbf{Partial credit:} None
% \textbf{Trick answer:} A (negative covariance/correlation suggests orthogonality, but $|\rho| > 0.5$ means angles are too extreme)
% \textbf{Trick answer:} B (uses covariance magnitude instead of correlation; $C_{12} = 3$ and $C_{34} = 1.2$ are small but this doesn't directly indicate angles)
% \textbf{Wrong:} C (identifies one correct pair but misses the others)
% \textbf{Trick answer:} D (assumes negative correlation means near-orthogonal, but need $|\rho| < 0.5$, not just $\rho < 0$)

\divider

%%%%%%%%%%%%%%%%%%%%%%%%%%%%%%%%%%%%%%%%%%%%%%%%%%%%%%%%%%%%%%%%%
\newpage
%%%%%%%%%%%%%%%%%%%%%%%%%%%%%%%%%%%%%%%%%%%%%%%%%%%%%%%%%%%%%%%%%



{\bf PROBLEM 18:}
% WEEKS = 10, 11 SYNTHESIS
% TOPICS = SVD geometry, PCA variance structure, A^T A vs covariance

A centered data matrix $X \in \mathbb{R}^{100 \times 3}$ represents 100 observations of three variables. The matrix $X^TX$ has eigenvalues:
$$\lambda_1 = 1600, \quad \lambda_2 = 400, \quad \lambda_3 = 100$$

Consider the linear transformation $T: \mathbb{R}^3 \to \mathbb{R}^{100}$ defined by $T(\mathbf{v}) = X\mathbf{v}$. This transformation maps the unit sphere in $\mathbb{R}^3$ to an ellipsoid in $\mathbb{R}^{100}$.

Which statement correctly relates the SVD geometry to the PCA variance structure?

\begin{enumerate}[(A)]
\item The ellipsoid has semi-axis lengths 1600, 400, and 100, which equal the variances along the three principal component directions.
\item The ellipsoid has semi-axis lengths 40, 20, and 10, which equal the variances along the three principal component directions.
\item The ellipsoid semi-axes have lengths in ratio 4:2:1, while the principal component variances have ratio 16:4:1.
\item The ellipsoid semi-axes have lengths in ratio 4:2:1, and the principal component variances have the same ratio 4:2:1.
\item The ellipsoid lies in a 3-dimensional subspace of $\mathbb{R}^{100}$, but the semi-axis lengths and PC variances are unrelated quantities.
\end{enumerate}

% \textbf{Correct Answer:} C
% \textbf{Explanation:} This requires connecting SVD geometry to PCA statistics:
% 
% 1) SVD GEOMETRY: The eigenvalues of $X^TX$ are $\sigma_i^2$, so $\sigma_1 = 40$, $\sigma_2 = 20$, $\sigma_3 = 10$. 
%    These singular values are the semi-axis lengths of the ellipsoid.
%    Ratio: $40 : 20 : 10 = 4 : 2 : 1$.
% 
% 2) PCA VARIANCE: The covariance matrix is $C = X^TX/(n-1) = X^TX/99$, with eigenvalues 
%    $\lambda_1^{(C)} \approx 16.16$, $\lambda_2^{(C)} \approx 4.04$, $\lambda_3^{(C)} \approx 1.01$.
%    These are the PC variances. Ratio: approximately $16 : 4 : 1$.
% 
% 3) KEY INSIGHT: Semi-axis lengths scale as $\sigma_i$ (ratio 4:2:1), while variances scale as 
%    $\sigma_i^2/(n-1)$ (ratio 16:4:1). The linear vs. quadratic relationship explains the difference.
%    Both capture the same ellipsoidal structure: singular values measure stretching factors (distances), 
%    while variances measure spread (squared distances, normalized by sample size).
% 
% \textbf{Partial credit:} D (correctly identifies semi-axis ratio but misses that variance ratio is the square)
% \textbf{Trick answer:} A (confuses eigenvalues of $X^TX$ with both semi-axes and variances; eigenvalues are $\sigma_i^2$, not the lengths or variances directly)
% \textbf{Trick answer:} B (correctly computes semi-axis lengths as $\sqrt{\lambda_i}$, but these are NOT the variances; variances require division by $n-1$)
% \textbf{Wrong:} E (the quantities are deeply related: variance $\propto \sigma_i^2$, so structure is preserved)

\divider




{\bf PROBLEM 19:}
% WEEK = 10
% TOPICS = Frobenius norm from singular values

A matrix $A \in \mathbb{R}^{4 \times 3}$ has singular values $\sigma_1 = 6$, $\sigma_2 = 3$, $\sigma_3 = 2$.

What is the Frobenius norm $\|A\|_F$?

\begin{enumerate}[(A)]
\item $11$
\item $49$
\item $7$
\item $\sqrt{11}$
\item Not enough information to determine $\|A\|_F$.
\end{enumerate}

% \textbf{Correct Answer:} C
% \textbf{Explanation:} The Frobenius norm can be computed directly from singular values:
% $$\|A\|_F = \sqrt{\sigma_1^2 + \sigma_2^2 + \sigma_3^2} = \sqrt{36 + 9 + 4} = \sqrt{49} = 7$$
% This is equivalent to $\sqrt{\sum_{i,j} a_{ij}^2}$ (the entry-wise definition), but the singular value formula shows that $\|A\|_F$ depends only on the singular values, not on the specific entries or the orthogonal factors $U$ and $V$.
% \textbf{Partial credit:} None
% \textbf{Trick answer:} A (adds singular values $6 + 3 + 2 = 11$ instead of computing $\sqrt{\sum \sigma_i^2}$; confuses sum with norm)
% \textbf{Trick answer:} B (computes $\sigma_1^2 + \sigma_2^2 + \sigma_3^2 = 49$ but forgets the square root)
% \textbf{Trick answer:} D (takes square root of the sum of singular values $\sqrt{11}$; combines both errors)
% \textbf{Trick answer:} E (believes one needs the actual matrix entries or the orthogonal factors; doesn't recognize that singular values fully determine the Frobenius norm)

\divider


%%%%%%%%%%%%%%%%%%%%%%%%%%%%%%%%%%%%%%%%%%%%%%%%%%%%%%%%%%%%%%%%%
\newpage
%%%%%%%%%%%%%%%%%%%%%%%%%%%%%%%%%%%%%%%%%%%%%%%%%%%%%%%%%%%%%%%%%

{\bf PROBLEM 20:}
% WEEK = 10
% TOPICS = SVD structure, fundamental subspaces, left vs right singular vectors

The figure shows the full SVD $A = U\Sigma V^T$ for a matrix $A \in \mathbb{R}^{12 \times 5}$.
Filled circles represent nonzero singular values; empty circles represent zero singular values.
Vertical lines indicate columns; horizontal lines indicate rows.

\begin{center}
\begin{figure}[h]
\includegraphics[width=7in]{2030 Q5 SVD ICON.jpg}
%\caption{[left] for PROBLEM 14 ; [right] for PROBLEM 15.}
\end{figure}
\end{center}
Which statement is most TRUE?

\begin{enumerate}[(A)]
\item $\dim(\ker A) = 9$.
\item The first three columns of $V$ form an orthonormal basis for $\mathrm{im}(A)$.
\item The last 9 columns of $U$ form an orthonormal basis for the codomain of $A$.
\item The last two columns of $V$ form an orthonormal basis for $\ker(A)$.
\item The rank of $A$ equals $5$.
\end{enumerate}

% \textbf{Correct Answer:} D
% \textbf{Explanation:} For $A \in \mathbb{R}^{12 \times 5}$ with rank $r = 3$:
% - Right singular vectors $\mathbf{v}_i \in \mathbb{R}^5$ (domain)
% - Left singular vectors $\mathbf{u}_i \in \mathbb{R}^{12}$ (codomain)
% - $\{\mathbf{v}_1, \mathbf{v}_2, \mathbf{v}_3\}$ spans $\mathrm{im}(A^T)$ (row space)
% - $\{\mathbf{v}_4, \mathbf{v}_5\}$ spans $\ker(A)$ ← this is (D), TRUE
% - $\{\mathbf{u}_1, \mathbf{u}_2, \mathbf{u}_3\}$ spans $\mathrm{im}(A)$ (column space)
% - $\{\mathbf{u}_4, \ldots, \mathbf{u}_{12}\}$ spans $\ker(A^T)$ (left nullspace)
%
% \textbf{Partial credit:} None
% \textbf{Trick answer:} A (computes $12 - 3 = 9$, which is $\dim(\ker A^T)$, not $\dim(\ker A)$; 
%   confuses left nullspace with right nullspace. Correct: $\dim(\ker A) = 5 - 3 = 2$)
% \textbf{Trick answer:} B (right singular vectors span $\mathrm{im}(A^T)$, not $\mathrm{im}(A)$; 
%   the column space is spanned by LEFT singular vectors $\mathbf{u}_1, \mathbf{u}_2, \mathbf{u}_3$)
% \textbf{Trick answer:} C (the vectors $\mathbf{u}_i$ live in $\mathbb{R}^{12}$, not $\mathbb{R}^5$; 
%   confuses domain with codomain)
% \textbf{Trick answer:} E (rank = number of nonzero singular values = 3, not the size of $V$; 
%   a matrix can have full-size orthogonal factors while being rank-deficient)

\divider


\end{document}